% EXPECTED-LEAKS: IEEEauthorblockN, IEEEauthorblockA, IEEEkeywords, thanks
\documentclass[conference]{IEEEtran}
\usepackage{cite}
\usepackage{amsmath,amssymb}
\usepackage{graphicx}
\usepackage{hyperref}

\begin{document}
\title{A Scalable Approach to Distributed Consensus}

\author{\IEEEauthorblockN{Alice Smith\IEEEauthorrefmark{1} and Bob Jones\IEEEauthorrefmark{2}}
\IEEEauthorblockA{\IEEEauthorrefmark{1}Dept.\ of Computer Science, Example University\\
Email: alice@example.edu}
\IEEEauthorblockA{\IEEEauthorrefmark{2}Research Institute\\
Email: bob@example.org}}

\maketitle

\begin{abstract}
We present a new distributed consensus protocol that scales to thousands of nodes while maintaining linearisability. Experimental evaluation shows a 3x throughput improvement over prior work \cite{raft2014}.
\end{abstract}

\begin{IEEEkeywords}
consensus, distributed systems, scalability
\end{IEEEkeywords}

\section{Introduction}
Distributed consensus underpins modern fault-tolerant systems. Prior approaches such as Paxos \cite{paxos1998} and Raft \cite{raft2014} rely on a single leader, creating a throughput ceiling.

\section{Related Work}
Lamport's original Paxos paper \cite{paxos1998} introduced the classic single-decree formulation. Multi-Paxos \cite{paxos-live} later extended this to log replication.

\section{Protocol}
Our protocol operates in two phases. In phase one, nodes exchange proposals using formula $v = f(t, q)$. In phase two, commit happens when $|q| \geq n/2 + 1$.

\begin{equation}
\Pr[\text{agreement}] \geq 1 - e^{-n/c}
\label{eq:agreement}
\end{equation}

Equation~\eqref{eq:agreement} bounds the failure probability.

\section{Evaluation}
Figure~\ref{fig:tput} shows throughput as a function of cluster size.

\section{Conclusion}
Future work includes Byzantine fault tolerance.

\bibliographystyle{IEEEtran}
\bibliography{refs}

\end{document}
